\section{F. Statements usable in the later sections}

We use $\Gamma_{g,\nu}$ for the extended geometric group permitting puncture permutations and $\Gamma_{g,\nu}^{!+}$ for the pure oriented group, following \S19, folio 102, scans 204--205. In each formula below the chosen variant is stated explicitly. The full profinite loop-outer group is $B=\Out_{\rm lac}(P)$; the closure $H=\overline M$ of a chosen geometric subgroup $M$ is a subgroup of $B$. We do not identify that closure with the abstract profinite completion without the congruence assertion required for that identification.

\begin{claim}{F.1}{Proved}{LM4-F.a: fixed discretizations and prediscretizations}
\label{LM4-F.a}
For $M=\Gamma_{g,\nu}$ and $H=\overline M\leq B$, the sets of outer discretizations and prediscretizations are $B/M$ and $B/H$. For a given finite $G\leq B$ their fixed sets are exactly D.16's transporter sets and centralizer-orbit decompositions. In the three types of D.8, every fixed prediscretization has a fixed outer discretization above it. The same lifting conclusion holds for a finite solvable $G$ and genus at most two for the orientation-preserving geometric group covered by Boggi's Theorem C. Over a fixed coset represented by $b$, the fibre is
\[
 (H/M)^{b^{-1}Gb},
\]
with its full labelled-embedding decomposition as in D.16.
\end{claim}
\begin{proof}
The identifications and the fixed-set computation are proved in D.16, including their inner-automorphism ambiguities. A fixed coset $bH$ has $K=b^{-1}Gb\leq H$. Its fibre consists of $bhM$, $h\in H$; this coset is fixed precisely when $h^{-1}Kh\leq M$, proving the formula. In the three low types D.8 conjugates every finite $K\leq H$ into $M$, giving a point in that fibre. In genus at most two and for solvable $K$, Boggi's Theorem C \cite{R-Boggi} gives the same conjugation into the oriented geometric group. The congruence subgroup property in these genera identifies its completion with $H$, so the theorem applies to the displayed subgroup. We retain the labelled embeddings in the fibre decomposition because subgroup conjugacy alone does not fix an automorphism of $K$.

The actual based version has $A(P)/J$ in place of $B/M$, as D.16 proves. A chosen genuine lift of $G$ acts on it; an outer subgroup alone acts on its quotient by $P$. These are the exact maps through which the fixed-set assertions of \S28, folios 178--181, scans 352--358 are used. For a finite group in the full $B$ not known to stabilize a prediscretization, the same formula still determines the fixed set as a transporter; its nonemptiness is exactly the unresolved recognition question D.12, not an additional conclusion of this lemma.
\end{proof}

\begin{claim}{F.2}{Refuted}{LM4-F.b, finiteness and equality of the two centralizers}
\label{LM4-F.b}
The centralizer of a finite subgroup $G\leq\Gamma^+_{g,\nu}$ need not be finite, either discretely or in the geometric completion. In genus at most two, for finite solvable $G$, the normalizer comparison and the open-kernel argument below give
\[
 C_{\widehat{\Gamma^+_{g,\nu}}}(G)
       =\overline{C_{\Gamma^+_{g,\nu}}(G)}.
\]
The corresponding equality with the full loop-outer group in place of the geometric completion is false already for $(g,\nu)=(0,3)$ and $G=S_3$.
\end{claim}
\begin{proof}
On a closed genus-two surface let $G$ be generated by the hyperelliptic involution. This involution is central in $\Gamma_2^+$, so its centralizer is the entire group. A Dehn twist about a nonseparating curve has infinite order: its action on integral homology sends a transverse vector $b$ to $b+a$, and its $n$th power sends it to $b+na$. This also gives infinite order in the congruence image, since some reduction modulo an integer detects every nonzero $n$. Thus the completed centralizer is infinite as well. For finite solvable $G$, the normalizer comparison, with the source induction and its topology input audited in F.7, gives $N_{\widehat M}(G)=\overline{N_M(G)}$ for $M=\Gamma^+_{g,\nu}$. The continuous conjugation homomorphism from this profinite normalizer to the finite group $\Aut(G)$ has open kernel $C_{\widehat M}(G)$. Intersecting the dense subgroup $N_M(G)$ with that open kernel gives a dense subgroup of the kernel: every open subset of the kernel is open in the normalizer and meets $N_M(G)$. The intersection is exactly $C_M(G)$, proving the displayed centralizer equality. The map of the completion to the loop-outer group is injective in these genera by the congruence subgroup property. The cited sources state the equality for every finite $G$; F.7 records the precise missing intersection step for their nonsolvable assertion. The passage from any such equality to finiteness would still require the discrete centralizer to be finite, contradicted by the hyperelliptic example.

For the smallest hyperbolic type, $\Gamma_{0,3}^+=S_3$, and $C_{S_3}(S_3)=1$. Harbater--Schneps identify $C_B(S_3)=GT_0$. This group is infinite: its multiplier contains the cyclotomic image $\widehat\Z^\times$ of $G_\Q$. These are actual homomorphisms on the same group $\widehat F_2$, with the geometric $S_3$ supplied by rational M\"obius transformations, so the commutation with the Galois image is the naturality of the fundamental-group action. The two centralizers are therefore $1$ and an infinite group. With orientation reversal included the geometric centralizer becomes $C_2$, and the same failure remains.

The statement denoted (c) on scans 383--385, \S30, folios 194--195, concerns the centralizer of an arithmetic germ $\Gamma^{\natural}$. It cannot be replaced merely by a finite $G$: scan 385 itself discusses infinite centralizers of finite groups. The counterexamples above compute the result of the requested replacement explicitly and explain which datum that replacement removes.
\end{proof}

\begin{claim}{F.3}{Proved}{LM4-F.b, the intrinsic open-image statement}
Let $M\leq B$ be an oriented geometric subgroup, $H=\overline M$, $N=N_B(H)$, and let $i:G\hookrightarrow H$ be a fixed finite embedding. Conjugation defines an action of $N/H$ on the set of $H$-conjugacy classes of embeddings $G\hookrightarrow H$. The image of the continuous homomorphism
\[
 C_N(i(G))\longrightarrow N/H
\]
is exactly the stabilizer of $[i]$, and its kernel is $C_H(i(G))$. When the orbit of $[i]$ is finite, this image is open. The orbit is finite for solvable finite $G$ in genus at most two and for every finite $G$ in the three low types of D.8.
\end{claim}
\begin{proof}
An element $nH$ stabilizes $[i]$ exactly when there is $h\in H$ with
\[
 ni(g)n^{-1}=hi(g)h^{-1}\quad(g\in G).
\]
This is equivalent to $h^{-1}n\in C_N(i(G))$. It proves the image assertion, and the kernel is obtained by intersecting the centralizer with $H$. Because $H$ is compact, the stabilizer is the image of a compact closed transporter and is closed in $N/H$. A closed subgroup of finite index in a profinite group is open: its finitely many closed cosets have open complements, so each coset is open.

For the claimed cases, finite subgroup conjugacy classes in $H$ come from the finite list of geometric classes, by D.8 or Boggi's solvable statement. There are finitely many geometric classes for each fixed finite group: a finite action has an orbifold quotient of bounded signature by the Euler characteristic formula; there are finitely many epimorphisms from each finitely generated orbifold group to the fixed finite group. Passing from subgroup classes to labelled-embedding classes adds at most $|\Aut(G)|$ choices per subgroup class. Thus the orbit is finite and the image is open. F.6 constructs the arithmetic quotient directly from universal monodromy and proves the open-image assertion in that quotient for every finite geometric subgroup and every hyperbolic type. The inclusion of that quotient into $N/H$ is retained there as the exact comparison map.
\end{proof}

\begin{claim}{F.4}{Proved}{LM4-F.c: the exact canonical lift of a section stabilizer}
\label{LM4-F.c}
Let $\pi$ be a centre-free discrete surface group or its full profinite completion, and put
\[
 A=\Aut_{\rm lac}(\pi),\qquad B=\Out_{\rm lac}(\pi).
\]
Let $G\leq B$ be finite and let $s:G\to A$ be a specified lift with $C_\pi(s(G))=1$. The natural map
\[
 C_A(s(G))\longrightarrow C_B(G)
\]
is injective. Its image is the stabilizer of the $\pi$-conjugacy class $[s]$ in the set of lifts of the labelled inclusion $G\hookrightarrow B$. Its inverse is the canonical homomorphism lifting that stabilizer to $A$ and centralizing $s(G)$. In the profinite case it is continuous. For geometrically realized actions with a nontrivial rotation the section-class set is finite, so this stabilizer has finite index and is open in the profinite case.
\end{claim}
\begin{proof}
The kernel is the inner automorphisms centralizing $s(G)$, identified, by centre-freeness, with $C_\pi(s(G))$. For $b\in C_B(G)$ choose an automorphism representative $v$. The conjugate homomorphism $g\mapsto vs(g)v^{-1}$ still lifts the labelled inclusion, because $b$ centralizes each $g$. It represents $[s]$ exactly when some $p\in\pi$ makes $\Inn(p)^{-1}v$ commute with every $s(g)$. This proves that the image equals the claimed stabilizer. The kernel computation proves uniqueness. The inverse is a homomorphism by uniqueness, since products of commuting lifts still commute with $s(G)$.

For a finitely generated profinite $\pi$, $A$ and $B$ are profinite, and the centralizer is closed. The map is a continuous injection from a compact space into a Hausdorff space, so its inverse on its image is continuous. In the discrete geometric situation Part B proves finiteness of the lift classes; D.3 proves their equality with the profinite classes. The orbit of a class is therefore finite. Its stabilizer is closed, as the image of the compact centralizer, hence open. This proves all topology assertions as well as the group assertion.

For the notation of scans 390--393, \S31, folios 197--199, the target $\mathfrak S(\pi_0)^G$ is the centralizer of the chosen genuine lift inside the loop automorphism group. The source is the stabilizer of the corresponding section class inside $\mathcal T(\pi_0)^G=C_{\Out_{\rm lac}(\pi_0)}(G)$. A point $P$ determines that section class after a base path; changing the path conjugates the resulting lift by $\pi_0$, exactly as displayed. Thus the canonicity belongs to the chosen section, and the geometric point supplies its conjugacy class.
\end{proof}

\begin{claim}{F.5}{Refuted}{The unrestricted uniqueness assertion for reflection ovals}
For a reflection fixing an oval, the lifting map in F.4 has kernel $\Z$ in the discrete case and $\widehat\Z$ after completion. Consequently it is not injective and its lifts are not unique without further data.
\end{claim}
\begin{proof}
Part B and D.4 identify $C_\pi(s(G))$ with the primitive oval subgroup and its closure. Every element of that subgroup induces an inner automorphism centralizing the chosen reflection lift and maps to the identity outer class. The kernel computation in F.4 therefore gives the displayed groups. A nonidentity primitive oval element is an explicit second lift of the identity outer element. The corrected general statement is the exact sequence with this kernel and with image the same section stabilizer, proved by the same representative-and-conjugator calculation. No splitting of that sequence is supplied by the point class alone.
\end{proof}
